%%═══════════════════════════════════════════════════════════════════
%% Lecture 03 — Band Matrices, Least Squares, Gram-Schmidt, and Eigenvalues
%% 77315 — Computational Physics
%% Date: 23/04/2026
%%═══════════════════════════════════════════════════════════════════

\section{Lecture 03 — Band Matrices, Least Squares, Gram-Schmidt, and
  Eigenvalues}\label{sec:lec03}

\begin{center}
\begin{tabular}{ll}
  \textbf{Date:}\quad 23/04/2026 & \\
\end{tabular}
\end{center}
\hrule\vspace{1.5em}

%%──────────────────────────────────────────────────────────────────
\subsection{Large and Sparse Matrices (Band Matrices)}
In real physical systems, ``everything does not interact with everything.''
Consider the 1D heat diffusion equation on a wire discretised into 10,000
segments. The temperature at segment $i$ only directly influences its
immediate neighbours $i-1$ and $i+1$. The interaction matrix $A$ is
$10{,}000 \times 10{,}000$ (100 million elements), but $99.99\%$ are zero —
only the main diagonal and two adjacent diagonals are non-zero. This is a
\textbf{Band Matrix}.

By storing and operating only on the band of width $n_b$, Gaussian
elimination time drops from $\bigO{n^3}$ to $\bigO{n_b^2\, n}$ and memory
from gigabytes to megabytes.

%%──────────────────────────────────────────────────────────────────
\subsection{Least Squares}
When fitting a line to noisy experimental data, we have an overdetermined
system $Ax = b$ with more equations (data points) than unknowns. A perfect
solution is impossible.

\noindent\textbf{Physical intuition:} The columns of $A$ span the ``model
space''. The data vector $b$ lives outside this space due to noise. The best
parameters $x$ are found by dropping a perpendicular from $b$ onto the model
space. The error vector $(b - Ax)$ must be orthogonal to every column of $A$:
\[
A^T(b - Ax) = 0 \implies \boxed{A^T A\, x = A^T b}
\]
These are the \textbf{Normal Equations}. Solving them minimises the squared
error $\norm{Ax - b}^2$.

%%──────────────────────────────────────────────────────────────────
\subsection{Gram-Schmidt Orthogonalization}
Explicitly forming $A^T A$ is numerically dangerous because it squares the
condition number of the matrix, violently amplifying floating-point
inaccuracies. Instead, \textbf{Gram-Schmidt} orthogonalises the columns of
$A$ directly.

\noindent\textbf{Physical intuition:} For each vector $a_j$, subtract its
projection onto all previously established orthogonal basis vectors:
\begin{equation}
a'_j = a_j - \frac{a_1 \cdot a_j}{a_1 \cdot a_1}\, a_1
\end{equation}
Repeating for all $j$ yields mutually orthogonal columns.

\textbf{Connection to QR decomposition.} Gram-Schmidt is algebraically
equivalent to a QR factorisation $A = \hat{A}R$, where $\hat{A}$ has
orthonormal columns ($\hat{A}^T\hat{A} = I$). Substituting into the Normal
Equations:
\begin{equation}
A^T A\, x = A^T b \;\implies\; Rx = \hat{A}^T b
\end{equation}
This upper-triangular system is solved by back-substitution in $\bigO{m^2}$,
entirely bypassing formation of $A^T A$.

\begin{tcolorbox}[colback=orange!5!white,colframe=orange!75!black,
  title=Numerical Stability Warning: Loss of Orthogonality]
If the columns of $A$ are nearly linearly dependent, the remainder after
subtraction can be many orders of magnitude smaller than the original. Once
intermediate values approach machine epsilon ($\varepsilon_{\text{mach}}
\approx 10^{-15}$), the subtraction is dominated by floating-point noise and
\textbf{orthogonality is silently lost}. In practice, use
\emph{Modified Gram-Schmidt} (subtracting projections one at a time) or
Householder reflections.
\end{tcolorbox}

%%──────────────────────────────────────────────────────────────────
\subsection{Introduction to Eigenvalues --- Jacobi Method}
For a real symmetric matrix $A$, we seek $\lambda_i$ such that
$A\varphi_i = \lambda_i\varphi_i$. The \textbf{Jacobi Method} applies a
sequence of plane-rotation transformations $T^T A T$, where each $T$ is a
Givens rotation that zeros one off-diagonal pair $(A_{pq}, A_{qp})$.

\noindent\textbf{Physical intuition:} At each step, take the 2D slice of
the $N$-dimensional matrix containing the largest off-diagonal element and
physically rotate the coordinate system in that plane until the element
vanishes. Convergence is guaranteed by two invariants:
\begin{enumerate}
    \item \textbf{Trace Invariance}: $\mathrm{Tr}(T^T A T) = \mathrm{Tr}(A)$.
    The sum of eigenvalues never changes.
    \item \textbf{Frobenius Norm Invariance}: $\sum_{i,j}(T^T A T)_{ij}^2 =
    \sum_{i,j} A_{ij}^2$. Total squared weight is conserved under orthonormal
    transformations.
\end{enumerate}

\begin{tcolorbox}[colback=green!5!white,colframe=green!50!black,
  title=Why the Algorithm Must Converge]
The Frobenius norm is shared between diagonal and off-diagonal entries. Since
the total is fixed, weight removed from off-diagonal must go to the diagonal.
After each full cycle of $n(n-1)/2$ rotations:
\begin{equation}
S_{\text{off}}^{(k+1)} \leq S_{\text{off}}^{(k)} \cdot e^{-1}
\end{equation}
so $S_{\text{off}}$ decreases by at least $e^{-1}$ per cycle and convergence
to a diagonal matrix is guaranteed.
\end{tcolorbox}

%%──────────────────────────────────────────────────────────────────
\subsection*{Recommended Reading}
\begin{itemize}
    \item \textit{Numerical Recipes}, W.H. Press et al.\ --- Ch.\ 2.4 (Band
    Diagonal Systems), Ch.\ 15.4 (General Linear Least Squares), Ch.\ 11.0
    (Eigensystems Introduction).
    \item \textit{An Introduction to Computational Physics}, Tao Pang ---
    Ch.\ 4 (Matrix Operations).
\end{itemize}
